% Generated draft: requires mathematical and source review.
\documentclass[11pt]{article}
\usepackage[margin=25mm]{geometry}
\usepackage{fontspec}
\setmainfont{Latin Modern Roman}
\usepackage{amsmath,amssymb,mathtools,unicode-math}
\setmathfont{Latin Modern Math}
\usepackage{enumitem,needspace,xurl,hyperref}
\setlength{\parindent}{0pt}
\setlength{\parskip}{5pt}
\setlength{\emergencystretch}{3em}
\newcommand{\N}{\mathbb{N}}
\newcommand{\Z}{\mathbb{Z}}
\newcommand{\Q}{\mathbb{Q}}
\newcommand{\R}{\mathbb{R}}
\newcommand{\C}{\mathbb{C}}
\newcommand{\F}{\mathbb{F}}
\newcommand{\A}{\mathbb{A}}
\newcommand{\PP}{\mathbb{P}}
\newcommand{\E}{\mathbb{E}}
\newcommand{\eps}{\varepsilon}
\newcommand{\dd}{\,\mathrm{d}}
\begin{document}
\section*{UnsolvedMath ID: 30004216}
\textbf{Source problem code:} catalogue entry
\subsection{Definitions and mathematical statement}
For real parameters \(t_1<\cdots<t_n\), the usual cyclic-polytope construction is
\[
C_{n,d}=\operatorname{conv}\{(t_i,t_i^2,\ldots,t_i^d):1\le i\le n\}\subseteq\R^d.
\]
The supplied record does not specify the admissible ranges of \(n,d\), or whether a particular realization is intended.

For a poset \(P\), let \(\mathcal O(P)\) be its order complex: its vertices are the elements of \(P\), and its nonempty simplices are the finite nonempty chains in \(P\). Write \(|\mathcal O(P)|\) for its geometric realization. For an integer \(m\ge0\), let
\[
S^m=\{x\in\R^{m+1}:\|x\|_2=1\}.
\]

\textbf{Source question, with unresolved specifications made explicit.} Let \(\mathcal P_{\mathrm{hyp}}(C_{n,d};k)\) denote the poset intended by the source's phrase ``hypersimplicial subdivisions.'' This notation is a placeholder: the supplied record does not define its elements or order relation. For every admissible triple \((n,d,k)\) intended by the source, is
\[
\bigl|\mathcal O\bigl(\mathcal P_{\mathrm{hyp}}(C_{n,d};k)\bigr)\bigr|
\simeq S^{\,n-d-2},
\]
where \(\simeq\) denotes homotopy equivalence?

A fully specified conjecture cannot be recovered from this record alone. In particular, the domain and range of \(k\), the construction defining hypersimplicial subdivisions, the ordering, and any restriction to proper subdivisions are missing. Whether an indiscrete subdivision is omitted is material: a poset containing a greatest or least element has contractible order complex. The record also gives no convention for a negative sphere dimension, including \(S^{-1}\) when \(n=d+1\). These qualifications remain unresolved rather than being imposed as additional hypotheses.
\subsection{Short English statement}
The question asks whether the order complex of the intended poset of hypersimplicial subdivisions associated with \(C_{n,d}\) and \(k\) is homotopy equivalent to \(S^{n-d-2}\). The supplied record leaves the poset, parameter ranges, and boundary conventions unspecified, so this is a statement-only draft requiring clarification and human mathematical review.
\subsection{Sources}
\begin{itemize}
\item Local UnsolvedMath registry, supplied complete record for permanent ID \(30004216\), titled ``Generalized Baues Problem for Hypersimplicial Subdivisions,'' in the Oberwolfach Reports Open Problems collection. Exact source problem code: \texttt{catalogue entry}. The record gives proposed year \(2019\) and identifies no proposer.
\item Supplied report citation: \emph{Geometric, Algebraic, and Topological Combinatorics} (2020), Oberwolfach Reports, pp.~2395--2472, DOI \texttt{10.4171/owr/2019/39}. \url{https://doi.org/10.4171/owr/2019/39}. The supplied background describes Oberwolfach Reports as documenting mathematical workshops and their open-problem discussions; no precise problem location within the cited page range is supplied.
\item Supplied literature reference: Olarte and Santos, ``hypersimplicial subdivisions.'' \url{https://doi.org/10.1007/s00029-021-00722-6}.
\item Supplied literature reference: Rambau and Santos, ``generalized Baues problem.'' \url{https://doi.org/10.1006/eujc.1999.0321}.
\item Supplied literature assessment, checked on \texttt{2026-08-22}: the generalized Baues question is reported solved for convex polygons and for \(k=1\), while the full cyclic-polytope hypersimplicial problem is reported to remain open. The supplied status is \texttt{partially\_solved}. This is supplied dated curation derived from the OWR record, not an independently verified literature assessment; the scope of the stated solved cases and the time-sensitive status require independent checking.
\item The registry labels the statement \texttt{corrected\_verified}, explaining that notation, typography, and sentence structure were repaired while mathematical content was preserved. This supplied editorial label does not establish independent mathematical verification of this draft or resolve its missing definitions.
\end{itemize}
\end{document}
